This section details our experimental methodology, including data preparation, fine-tuning protocols, and the evaluation framework used to benchmark ECAPA-TDNN and RawNet3 under resource-constrained conditions.

\subsection{Data Partitioning and Preprocessing}

To ensure strict zero-speaker leakage across splits, both datasets are partitioned into disjoint train, validation, and test subsets:

\begin{itemize}
    \item \textbf{TidyVoice:} The official training split is preserved for training (262{,}268 utterances, 3{,}666 speakers). The source development set (808 speakers) is partitioned 50/50 into Validation (29{,}720 utterances, 404 speakers) and Test (29{,}723 utterances, 404 speakers) using a deterministic seed (seed 42), optimizing for balanced utterance counts and language distribution with zero speaker overlap.
    \item \textbf{ViMD:} The training split contains 15{,}023 utterances from 10{,}291 speakers. The original test split (2{,}026 utterances, 1{,}344 speakers) is kept intact. A data audit identified two overlapping singleton speakers between validation and test (\texttt{spk\_73\_0186} and \texttt{spk\_76\_0219}); these singleton rows were pruned from the validation set to eliminate speaker leakage, yielding a clean validation split of 1{,}898 utterances across 1{,}318 speakers without reducing the available genuine trial capacity.
\end{itemize}

Table~\ref{tab:dataset-splits} summarizes the standardized data splits across both corpora.

\begin{table}[H]
\centering
\begin{tabular}{llrrr}
\toprule
\textbf{Dataset} & \textbf{Split} & \textbf{Utterances} & \textbf{Speakers} & \textbf{Genuine Pairs} \\
\midrule
\multirow{3}{*}{TidyVoice} 
  & Train      & 262{,}268 & 3{,}666 & ---   \\
  & Validation & 29{,}720  & 404     & 7{,}954 \\
  & Test       & 29{,}723  & 404     & 7{,}898 \\
\midrule
\multirow{3}{*}{ViMD} 
  & Train      & 15{,}023  & 10{,}291 & 7{,}044 \\
  & Validation & 1{,}898   & 1{,}318  & 863     \\
  & Test       & 2{,}026   & 1{,}344  & 1{,}042 \\
\bottomrule
\end{tabular}
\caption{Summary of standardized dataset splits and trial capacities.}
\label{tab:dataset-splits}
\end{table}

\textbf{Audio Preprocessing Pipeline.} All input audio files are standardized through a unified preprocessing pipeline:
\begin{enumerate}
    \item Decode raw audio to 32-bit floating-point arrays.
    \item Downmix multi-channel audio to mono via arithmetic averaging across channels.
    \item Resample audio to 16~kHz using polyphase filtering.
    \item Retain native acoustic dynamic range without amplitude normalization.
    \item Deterministically crop utterances: for training, audio is cropped to 48{,}000 samples (3.00~s) for ECAPA-TDNN and 48{,}240 samples (3.015~s) for RawNet3; for evaluation, segments are cropped to 48{,}000 samples for ECAPA-TDNN and 64{,}240 samples (4.015~s) for RawNet3.
    \item Utterances shorter than the target segment duration are extended via sample repetition rather than zero-padding to prevent introducing synthetic silence artifacts.
\end{enumerate}

\subsection{Fine-Tuning Configuration}

Both models are initialized from publicly released, state-of-the-art checkpoints pretrained on VoxCeleb:
\begin{itemize}
    \item \textbf{ECAPA-TDNN:} Loaded from \texttt{speechbrain/spkrec-ecapa-voxceleb} (revision \texttt{0f99f2d}), comprising 20.77M encoder parameters and outputting 192-dimensional embeddings.
    \item \textbf{RawNet3:} Loaded from \texttt{jungjee/RawNet3} (revision \texttt{c89102e}), comprising 16.28M encoder parameters and outputting 256-dimensional embeddings.
\end{itemize}

For target-domain adaptation, the original pretrained classification heads are discarded and replaced with freshly initialized Additive Angular Margin Softmax (AAM-Softmax) \cite{deng2019arcface} heads, with output dimensions matching the number of training speakers in each respective dataset (3{,}666 for TidyVoice, 10{,}291 for ViMD). During downstream verification and identification inference, the classification heads are removed and only the L2-normalized encoder embeddings are utilized.

\begin{table}[H]
\centering
\footnotesize
\renewcommand{\arraystretch}{1.0}
\begin{tabular}{lll}
\toprule
\textbf{Configuration} & \textbf{ECAPA-TDNN} & \textbf{RawNet3} \\
\midrule
Random seed           & 42                  & 42 \\
Objective function    & AAM-Softmax         & AAM-Softmax \\
AAM margin ($m$) / scale ($s$) & 0.2 / 30   & 0.2 / 30 \\
Precision             & FP16 (scale: 1024)  & FP16 (scale: 1024) \\
Batch size            & 24                  & 24 \\
Optimizer             & Adam                & Adam \\
Encoder learning rate & $1\times 10^{-4}$   & $1\times 10^{-4}$ \\
Head learning rate    & $1\times 10^{-4}$   & $1\times 10^{-4}$ \\
Weight decay          & $2\times 10^{-6}$   & $5\times 10^{-5}$ \\
Gradient clipping     & Global norm 5.0     & Global norm 5.0 \\
Max training epochs   & 3                   & 3 \\
Early stopping metric & Validation EER      & Validation EER \\
Early stopping patience & 1 epoch (tie: minDCF) & 1 epoch (tie: minDCF) \\
\bottomrule
\end{tabular}
\caption{Hyperparameter settings and training configurations.}
\label{tab:training-hyperparams}
\end{table}

\textbf{Resource-Constrained Protocol.} To operate within the remaining project-time budget, training uses one deterministic utterance per speaker per epoch (3{,}666 examples/epoch for TidyVoice and 10{,}291 examples/epoch for ViMD). Memory calibration with the 10{,}291-class ViMD head tested batch sizes 4, 8, 16, 24, and 32; batch size 32 was the largest tested passing size, not a measured VRAM limit. The predefined rule selected a batch size no greater than 80\% of that value, resulting in batch size 24, which was then confirmed by a three-batch real-audio optimization gate. Each architecture's two dataset runs were executed concurrently in a Kaggle T4$\times$2 session, with each Dataset assigned to it's single GPU.

\subsection{Evaluation Protocol and Metrics}

Verification performance is evaluated on frozen trial lists identical across all models to ensure an unbiased comparison:

\begin{itemize}
    \item \textbf{Trial Generation:} To prevent speakers with large utterance counts from dominating the evaluation, genuine pairs are capped at a maximum of 20 pairs per speaker. Each evaluation split is paired with 100{,}000 randomly sampled impostor trials, providing statistical resolution down to a False Acceptance Rate (FAR) of 0.01\%.
    \item \textbf{Scoring:} The trial score between an enrollment embedding $\boldsymbol{e}_1$ and a test embedding $\boldsymbol{e}_2$ is computed as the cosine similarity:
    \begin{equation}
        S(\boldsymbol{e}_1, \boldsymbol{e}_2) = \frac{\boldsymbol{e}_1 \cdot \boldsymbol{e}_2}{\|\boldsymbol{e}_1\|_2 \, \|\boldsymbol{e}_2\|_2}
    \end{equation}
    \item \textbf{Primary Metrics:} We report the Equal Error Rate (\textbf{EER} in \%), where False Acceptance Rate equals False Rejection Rate ($\text{FAR} = \text{FRR}$), and the minimum Detection Cost Function (\textbf{minDCF}) parameterized with $P_{\text{target}} = 0.01$ and $C_{\text{fa}} = C_{\text{miss}} = 1.0$:
    \begin{equation}
        C_{\text{Det}} = C_{\text{miss}} \cdot P_{\text{target}} \cdot \text{FRR} + C_{\text{fa}} \cdot (1 - P_{\text{target}}) \cdot \text{FAR}
    \end{equation}
    \item \textbf{Security Operating Point:} A security-critical threshold is calibrated on the Validation set at $\text{FAR} = 0.1\%$ and frozen. This threshold is then evaluated exactly once on the Test set to report operating accuracy, FAR, FRR, and True Acceptance Rate ($\text{TAR} = 1 - \text{FRR}$). Additionally, we compute $\text{TAR@FAR}$ across $\{5\%, 1\%, 0.1\%, 0.01\%\}$.
\end{itemize}
